\begin{derivation}[从\eqnrefs{eq:muon-full-w-amplitude-dp11}到\eqnrefs{eq:fermi-amplitude-pl-dp11}]
正文\eqnrefs{eq:muon-full-w-amplitude-dp11} 仍保留中间 $W$ 的完整动量依赖。实验中的 $\mu$ 子衰变满足 $|q^2|\ll m_W^2c^2$。在这一控制参数下展开完整 $W$ 传播子，局域四费米相互作用出现在最低阶。

传播子分母的几何展开为
\begin{align*}
\frac{1}{q^2-m_W^2c^2}
&=-\frac{1}{m_W^2c^2}\frac{1}{1-q^2/(m_W^2c^2)}\\
&=-\frac{1}{m_W^2c^2}
\left[1+\frac{q^2}{m_W^2c^2}
+\mathcal O\!\left(\frac{q^4}{m_W^4c^4}\right)\right].
\end{align*}
纵向分子与轻费米子流收缩时带有外腿质量。以 $\mu$ 子流为例，$q=p-k$，
\begin{align*}
q_\alpha\bar u_{\nu_\mu}(k)\gamma^\alpha P_Lu_\mu(p)
&=\bar u_{\nu_\mu}(k)(\slashed p-\slashed k)P_Lu_\mu(p)\\
&=m_\mu c\,\bar u_{\nu_\mu}P_Ru_\mu
-m_{\nu_\mu}c\,\bar u_{\nu_\mu}P_Lu_\mu.
\end{align*}
电子一侧同理得到
\begin{equation*}
q_\beta\bar u_e\gamma^\beta P_Lv_{\bar\nu_e}
=m_{\mathrm e} c\,\bar u_eP_Lv_{\bar\nu_e}
-m_{\nu_e}c\,\bar u_eP_Rv_{\bar\nu_e}.
\end{equation*}
因此纵向部分至少带两个轻费米子质量插入；相对于横向电流--电流项，它在振幅中受 $m_fm_{f'}/m_W^2$ 型比值抑制。保留最低阶 $\eta_{\alpha\beta}$ 后，先只看两个顶角与传播子的标量系数：
\begin{align*}
\left(-\frac{\ii g}{\sqrt2}\right)^2
\frac{-\ii}{q^2-m_W^2c^2}
&=\frac{\ii g^2}{2(q^2-m_W^2c^2)}\\
&=-\ii\frac{g^2}{2m_W^2c^2}
\left[1+\frac{q^2}{m_W^2c^2}
+\mathcal O\!\left(\frac{q^4}{m_W^4c^4}\right)\right].
\end{align*}
于是完整振幅在最低阶成为
\begin{align*}
\mathcal M
={}&-\frac{g^2}{2m_W^2c^2}
\bigl[\bar u_{\nu_\mu}\gamma^\alpha P_Lu_\mu\bigr]
\bigl[\bar u_e\gamma_\alpha P_Lv_{\bar\nu_e}\bigr]\\
&+\mathcal O\!\left(
\frac{q^2}{m_W^2c^2},
\frac{m_fm_{f'}}{m_W^2}
\right).
\end{align*}
这正是正文\eqnrefs{eq:fermi-amplitude-pl-dp11}。
\end{derivation}



\begin{derivation}[从\eqnrefs{eq:fermi-amplitude-pl-dp11}到\eqnrefs{eq:muon-lifetime-tree-dp11}]
正文\eqnrefs{eq:fermi-amplitude-pl-dp11} 已把虚 $W$ 消去，只剩两条左手带电流。自旋求和由 Dirac 完备关系计算，三体相空间则分解成连续两次 Lorentz 协变二体相空间积分。

为避免在相空间中反复携带 $c$，把四动量改写为能量量纲
\begin{equation*}
P\equiv cp,\qquad P'\equiv cp',\qquad K\equiv ck,\qquad K'\equiv ck',
\end{equation*}
并记 $M\equiv m_\mu c^2$。利用 $P_L=(1-\gamma^5)/2$ 与正文\eqnrefs{eq:gf-matching-dp11}，四费米振幅可等价写成通常的 $V-A$ 归一化
\begin{equation*}
\mathcal M
=-\frac{G_F^{(E)}}{\sqrt2}
\bigl[\bar u(K)\gamma^\alpha(1-\gamma^5)u(P)\bigr]
\bigl[\bar u(P')\gamma_\alpha(1-\gamma^5)v(K')\bigr].
\end{equation*}
忽略电子和中微子质量，并对初态 $\mu$ 子两个自旋取平均。由
\[
\sum_su(P,s)\bar u(P,s)=\slashed P+M,
\qquad
\sum_su(K,s)\bar u(K,s)=\slashed K,
\qquad
\sum_sv(K',s)\bar v(K',s)=\slashed K'
\]
得到
\begin{align*}
\overline{|\mathcal M|^2}
&=\frac{[G_F^{(E)}]^2}{4}\,
L^{\alpha\beta}H_{\alpha\beta},\\
L^{\alpha\beta}
&\equiv
\Tr\!\left[
\slashed K\gamma^\alpha(1-\gamma^5)
(\slashed P+M)\gamma^\beta(1-\gamma^5)
\right],\\
H_{\alpha\beta}
&\equiv
\Tr\!\left[
\slashed P'\gamma_\alpha(1-\gamma^5)
\slashed K'\gamma_\beta(1-\gamma^5)
\right].
\end{align*}
这里前因子 $1/4$ 同时包含振幅平方中的 $1/2$ 和初态自旋平均的 $1/2$。

先说明 $L^{\alpha\beta}$ 中的 $M$ 项为什么消失。利用
$(1-\gamma^5)\gamma^\beta=\gamma^\beta(1+\gamma^5)$，
\begin{align*}
&M\Tr\!\left[
\slashed K\gamma^\alpha(1-\gamma^5)
\gamma^\beta(1-\gamma^5)
\right]\\
&\quad=M\Tr\!\left[
\slashed K\gamma^\alpha\gamma^\beta
(1+\gamma^5)(1-\gamma^5)
\right]=0.
\end{align*}
因此只需计算四个普通 $\gamma$ 矩阵的迹与含一个 $\gamma^5$ 的迹。采用
\begin{align*}
\Tr(\gamma^\rho\gamma^\alpha\gamma^\sigma\gamma^\beta)
&=4(\eta^{\rho\alpha}\eta^{\sigma\beta}
-\eta^{\rho\sigma}\eta^{\alpha\beta}
+\eta^{\rho\beta}\eta^{\alpha\sigma}),\\
\Tr(\gamma^\rho\gamma^\alpha\gamma^\sigma\gamma^\beta\gamma^5)
&=-4\ii\epsilon^{\rho\alpha\sigma\beta},
\end{align*}
并用 $(1-\gamma^5)^2=2(1-\gamma^5)$，逐项得到
\begin{align*}
L^{\alpha\beta}
&=8\left(
K^\alpha P^\beta+K^\beta P^\alpha
-\eta^{\alpha\beta}K\!\cdot P
+\ii\epsilon^{\alpha\beta\rho\sigma}K_\rho P_\sigma
\right),\\
H_{\alpha\beta}
&=8\left(
P'_\alpha K'_\beta+P'_\beta K'_\alpha
-\eta_{\alpha\beta}P'\!\cdot K'
+\ii\epsilon_{\alpha\beta\mu\nu}P'^\mu K'^\nu
\right).
\end{align*}
把每个张量分成对称部分 $S$ 与反对称部分 $A$。由于 $S^{\alpha\beta}A_{\alpha\beta}=0$，交叉项消失。对称部分缩并为
\begin{equation*}
S_L^{\alpha\beta}S_{H,\alpha\beta}
=2\left[(K\!\cdot P')(P\!\cdot K')
+(K\!\cdot K')(P\!\cdot P')\right],
\end{equation*}
而反对称部分用
\begin{equation*}
\epsilon^{\alpha\beta\rho\sigma}
\epsilon_{\alpha\beta\mu\nu}
=-2\left(\delta^\rho_{\mu}\delta^\sigma_{\nu}
-\delta^\rho_{\nu}\delta^\sigma_{\mu}\right)
\end{equation*}
得到
\begin{equation*}
A_L^{\alpha\beta}A_{H,\alpha\beta}
=2\left[(K\!\cdot P')(P\!\cdot K')
-(K\!\cdot K')(P\!\cdot P')\right].
\end{equation*}
两者相加后第二种标量积精确相消，于是
\begin{equation*}
L^{\alpha\beta}H_{\alpha\beta}
=256(K\!\cdot P')(P\!\cdot K'),
\end{equation*}
从而
\begin{equation*}
\boxed{
\overline{|\mathcal M|^2}
=64[G_F^{(E)}]^2(P\!\cdot K')(P'\!\cdot K).}
\end{equation*}
关键结构现在显式可见：$V-A$ 的反对称 $\epsilon$ 张量项恰好删除了另一种标量积配对，而不是给平方振幅额外增加一个独立项。

对两个不可观测的中微子积分。定义
\begin{equation*}
Q\equiv K+K'=P-P',\qquad s_\nu\equiv Q^2.
\end{equation*}
三体 Lorentz 不变相空间先插入恒等分解
\[
1=\int\frac{\dd s_\nu}{2\pi}
\int\dd^4Q\,(2\pi)^4
\delta^{(4)}(Q-K-K')\delta(Q^2-s_\nu)\Theta(Q^0),
\]
于是严格因子化为
\begin{equation*}
\dd\mathfrak P_3
=\frac{\dd s_\nu}{2\pi}
\dd\mathfrak P_2(P;P',Q)\,
\dd\mathfrak P_2(Q;K,K').
\end{equation*}
固定 $Q$ 后只需计算二阶张量积分
\begin{equation*}
I^{\alpha\beta}(Q)
\equiv\int\dd\mathfrak P_2(Q;K,K')K^\alpha K'^\beta.
\end{equation*}
积分后只剩 $Q^\mu$ 与 $\eta^{\mu\nu}$，故
\begin{equation*}
I^{\alpha\beta}=A s_\nu\eta^{\alpha\beta}+B Q^\alpha Q^\beta.
\end{equation*}
系数 $A,B$ 不需要猜。到 $Q$ 静止系，令 $\sqrt{s_\nu}>0$，则
\begin{equation*}
K^\mu=\frac{\sqrt{s_\nu}}2(1,\mathbf n),\qquad
K'^\mu=\frac{\sqrt{s_\nu}}2(1,-\mathbf n),\qquad
\dd\mathfrak P_2=\frac{\dd\Omega}{32\pi^2}.
\end{equation*}
因此
\begin{align*}
I^{00}
&=\frac{s_\nu}{4}\frac{4\pi}{32\pi^2}
=\frac{s_\nu}{32\pi},\\
I^{0i}&=0,\\
I^{ij}
&=-\frac{s_\nu}{4}\frac{1}{32\pi^2}
\int\dd\Omega\,n_in_j
=-\frac{s_\nu}{96\pi}\delta_{ij}.
\end{align*}
与协变分解逐分量比较即得
\begin{equation*}
A=\frac1{96\pi},\qquad B=\frac1{48\pi},
\end{equation*}
即
\begin{equation*}
I^{\alpha\beta}(Q)
=\frac{1}{96\pi}
\left(s_\nu\eta^{\alpha\beta}+2Q^\alpha Q^\beta\right).
\end{equation*}
所以
\begin{align*}
&\int\dd\mathfrak P_2(Q;K,K')
(P\!\cdot K')(P'\!\cdot K)\\
&\qquad=P_\alpha P'_\beta I^{\beta\alpha}
=\frac{1}{96\pi}
\left[s_\nu(P\!\cdot P')+2(P\!\cdot Q)(P'\!\cdot Q)\right].
\end{align*}

在 $\mu$ 子静止系 $P=(M,\mathbf0)$，令
$x=2E_e/M$。无质量电子满足 $P'^2=0$，故
\begin{align*}
P\!\cdot P'&=\frac{M^2x}{2},\\
s_\nu=(P-P')^2&=M^2(1-x),\\
P\!\cdot Q&=M^2\left(1-\frac{x}{2}\right),\\
P'\!\cdot Q&=\frac{M^2x}{2}.
\end{align*}
代入可得
\begin{equation*}
\int\dd\mathfrak P_2(Q;K,K')\overline{|\mathcal M|^2}
=\frac{[G_F^{(E)}]^2M^4}{3\pi}x(3-2x).
\end{equation*}
外层二体相空间由一般公式
\begin{equation*}
\dd\mathfrak P_2(P;P',Q)
=\frac{|\mathbf P'|}{16\pi^2M}\dd\Omega_e
\end{equation*}
以及 $|\mathbf P'|=E_e=Mx/2$ 化为
\begin{equation*}
\dd\mathfrak P_2(P;P',Q)
=\frac{x}{32\pi^2}\dd\Omega_e,
\qquad
|\dd s_\nu|=M^2\dd x.
\end{equation*}
于是
\begin{align*}
\frac{\dd\Gamma_\mu^{(0)}}{\dd x}
&=\frac{1}{2M\hbar}\frac1{2\pi}
\left[\frac{[G_F^{(E)}]^2M^4}{3\pi}x(3-2x)\right]
\left[\frac{x}{32\pi^2}\int\dd\Omega_e\right]M^2\\
&=\frac{[G_F^{(E)}]^2M^5}{96\pi^3\hbar}
 x^2(3-2x),
\end{align*}
其中 $\int\dd\Omega_e=4\pi$。这就是正文\eqnrefs{eq:muon-differential-tree-dp73}。

总宽度现在只剩单变量初等积分：
\begin{align*}
\Gamma_\mu^{(0)}
&=\frac{[G_F^{(E)}]^2M^5}{96\pi^3\hbar}
\int_0^1\dd x\,x^2(3-2x)\\
&=\frac{[G_F^{(E)}]^2M^5}{96\pi^3\hbar}
\left[x^3-\frac{x^4}{2}\right]_0^1\\
&=\frac{[G_F^{(E)}]^2(m_\mu c^2)^5}{192\pi^3\hbar},
\end{align*}
即正文\eqnrefs{eq:muon-lifetime-tree-dp11}。
\end{derivation}

